\section{Related Work}

\label{sec:related}

\subsection{Agent Frameworks}

Early agent frameworks including LangChain~\cite{langchain2022} and
AutoGPT~\cite{autogpt2023} established the pattern of wrapping LLM inference
with tool-use capabilities. These frameworks provide useful abstractions but do
not address compute provisioning, multi-channel identity, or observability as
first-class concerns. The OpenAI Assistants API~\cite{openai2023gpt4} provides
managed tool execution but ties the agent runtime to a single model provider,
creating the vendor lock-in problem described in Section~\ref{sec:intro}.

\subsection{Coding Agents}

SWE-bench~\cite{jimenez2024swebench} established a rigorous benchmark for
software engineering agents, catalyzing a wave of systems including
Devin~\cite{cognition2024devin}, Aider~\cite{aider2024}, and SWE-
agent~\cite{yang2024sweagent}. These systems demonstrate that LLMs can resolve
real GitHub issues when equipped with appropriate tools and scaffolding.
However, they treat the execution environment as a fixed background assumption
rather than a dynamically managed resource.

\subsection{Computer Use}

Anthropic's computer use capabilities~\cite{anthropic2024computer} and
OpenAI's Operator~\cite{openai2025operator} enable LLMs to interact with
graphical interfaces. The Hanzo Operative framework~\cite{hanzo2024operative}
extends this with formal action specifications and safety boundaries. These
systems occupy a single layer of the UAH architecture---the compute tier---
without integrating the surrounding infrastructure.

\subsection{Model Context Protocol}

The Model Context Protocol (MCP)~\cite{mcp2024} defines a standard interface
for tool discovery and invocation, analogous to USB for AI tools. MCP
standardizes the tool protocol layer but leaves all other layers unspecified.
The UAH adopts MCP as the tool protocol standard while adding six additional
layers above and below.

\subsection{Multi-Agent Systems}

Research on multi-agent coordination~\cite{park2023generative,
wu2023autogen} has demonstrated emergent capabilities from collections of
specialized agents. The UAH provides the routing, shared state, and
observability infrastructure that makes multi-agent coordination tractable at
production scale.

% ============================================================================
% 3. THE SEVEN-LAYER ARCHITECTURE
% ============================================================================
